\documentclass[runningheads]{llncs}

\usepackage{graphicx}
\usepackage{booktabs}
\usepackage{amsmath}
\usepackage{amssymb}
\usepackage{url}
\usepackage{enumitem}
\usepackage{xcolor}

\title{From OWL to EER to Natural Language:\\How Representation Formality Shapes LLM Reasoning\\over Conceptual Models}

\titlerunning{Representation Formality and LLM Reasoning}

\author{
Wolfgang Maass\inst{1} \and
Veda C. Storey\inst{2} \and
Stephen W. Liddle\inst{3}
}

\authorrunning{W. Maass et al.}

\institute{
German Research Center for Artificial Intelligence (DFKI), Saarbr\"{u}cken, Germany \\
\email{wolfgang.maass@dfki.de}
\and
Georgia State University, Atlanta, GA, USA \\
\email{vstorey@gsu.edu}
\and
Brigham Young University, Provo, UT, USA \\
\email{liddle@byu.edu}
}

\begin{document}

\maketitle

% ══════════════════════════════════════════════════════════════
\begin{abstract}
We present the first benchmark for evaluating LLM reasoning over conceptual models at three formalization levels. Using a single controlled domain (Northwind) represented as an OWL~2 ontology, an EER diagram with natural-language annotations, and a natural-language requirements document, we test whether the classical formal/semi-formal/informal tripartition \cite{mylopoulos1992conceptual} predicts LLM output quality. Two open-weight models (Mistral-7B and Llama-3.1-8B) execute 190 tasks spanning SPARQL generation, OCL constraint synthesis, and entity/relationship extraction, each repeated $K\!=\!5$ times (1,900 total records). Our central result overturns the intuitive \emph{formalization gradient}: the naive expectation that quality decreases with formalization is \emph{not} supported once each level is scored on a defensible metric. Formal SPARQL queries run without error 73--74\% of the time but return the \emph{correct entities} only 19--27\% of the time (projection-tolerant execution match); semi-formal OCL captures the intended constraint at 90--93\%; and informal entity extraction, scored against each task's own gold, reaches 95--96\%. Rather than a monotonic quality decline, formalization governs \emph{verifiability and failure mode}: formal representations make errors machine-detectable (most SPARQL that runs is silently wrong), whereas informal outputs score high on forgiving coverage metrics. A robust \emph{model--formalization interaction} nonetheless emerges: Llama leads on formal SPARQL (27\% vs.\ 19\%) and strict OCL synthesis (87\% vs.\ 27\%), while both models converge on informal extraction. Both understand constraint semantics at similar rates (90--93\% loose OCL) but differ in syntactic fidelity. These findings identify a concrete automation boundary in the CM pipeline---formal generation demands execution verification, semi-formal synthesis a syntax-normalization layer, informal extraction is already reliable---and show why cross-level quality claims require a unified expert rubric rather than per-level automated proxies.
\end{abstract}

\textbf{Keywords:} conceptual modeling, formalization levels, OWL ontology, EER diagrams, large language models, SEQUAL, automation boundary

% ══════════════════════════════════════════════════════════════
\section{Introduction}
\label{sec:intro}

Conceptual modeling has long recognized that representations exist along a formalization spectrum. Mylopoulos \cite{mylopoulos1992conceptual} distinguished \emph{formal}, \emph{semi-formal}, and \emph{informal} representations, arguing that the choice of formalization level determines what operations---verification, transformation, automated reasoning---can be reliably performed over a model. With the rapid adoption of large language models (LLMs) as reasoning engines over structured and unstructured inputs, this classical distinction acquires practical urgency: does the formalization level of a conceptual model predict how well an LLM can reason over it?

We operationalize this question using three CM artifacts recognized by the ER community, all describing the same domain:

\begin{description}[noitemsep,style=nextline]
  \item[High-formal: OWL~2 ontology.]
    A Northwind domain ontology with TBox axioms, cardinality restrictions, and a populated ABox. Tasks: generate SPARQL queries answering competency questions. Correctness verified by execution.
  \item[Semi-formal: EER diagram with NL annotations.]
    The same domain as a structured prose EER linearization, with business rules expressed as natural-language annotations rather than formal constraints. Tasks: synthesize OCL constraints, reason about participation and cardinality.
  \item[Low-formal: NL requirements document.]
    The same domain described in business-analyst prose. Tasks: identify entities and relationships, flag ambiguities, elicit constraints.
\end{description}

\noindent The single-domain design eliminates the confound that plagues cross-task formalization studies: observed quality differences cannot be attributed to domain difficulty because all three levels describe the same enterprise reality (a food distribution company with customers, orders, products, suppliers, employees, and territories).

This paper makes three contributions:
\begin{enumerate}[noitemsep]
  \item A public benchmark of 190 CM reasoning tasks spanning OWL/EER/NL over a single controlled domain, with execution-verified and expert-annotated gold standards.
  \item Empirical identification of the automation boundary in the CM pipeline: where 7--8B LLMs can assist reliably (constraint comprehension, in-scope entity extraction) and where they cannot without support (executable SPARQL, where only 19--27\% of generated queries are correct despite 73--74\% executing; syntactically standard OCL).
  \item A model--formalization interaction finding: different model architectures have complementary strengths across the formalization spectrum, suggesting that formalization level should inform model selection in CM tool design.
\end{enumerate}

% ══════════════════════════════════════════════════════════════
\section{Related Work and Theoretical Grounding}
\label{sec:related}

\subsection{Formalization in Conceptual Modeling}

The formal/semi-formal/informal tripartition originates with Mylopoulos \cite{mylopoulos1992conceptual}, who argued that formal representations enable automated reasoning while informal representations require human interpretation. Wand and Weber \cite{wand1990ontological} showed that a modeling grammar's representational completeness determines what inferences are valid---formal representations reduce ambiguity by constraining the space of valid interpretations.

\subsection{SEQUAL Framework}

We ground our analysis in the SEQUAL framework \cite{krogstie1995quality,krogstie2016quality}, which decomposes model quality into syntactic quality (conformance to grammar), semantic quality (correspondence between model and domain), and pragmatic quality (correct interpretation by the audience). We map our formalization metric $F(T)$ onto SEQUAL dimensions:

\begin{itemize}[noitemsep]
  \item \textbf{Syntactic constraint density $S(T) \rightarrow$ Syntactic quality $Q_s$.} OWL~2 has a W3C grammar checkable by parser; EER annotations are unconstrained; NL has no grammar.
  \item \textbf{Evaluation decidability $D(T) \rightarrow$ Pragmatic quality $Q_p$.} SPARQL correctness is machine-verifiable ($D=1$); OCL is partially checkable ($D=0.5$); NL requires expert judgment ($D=0$).
  \item \textbf{Output-space cardinality $C(T) \rightarrow$ Semantic quality $Q_{se}$ (inverse).} Low cardinality (few correct answers) corresponds to high semantic quality; high cardinality (many acceptable answers) corresponds to low semantic quality.
\end{itemize}

\subsection{Moody's Physics of Notations}

Moody \cite{moody2009physics} identified properties that make notations effective for human interpreters. Three principles extend to LLMs via a training-data argument: (1)~\emph{semiotic clarity}---OWL constructs have unambiguous semantics, producing a consistent training signal; (2)~\emph{semantic transparency}---formal notations reduce per-token disambiguation burden; (3)~\emph{complexity management}---modular ontology structure reduces context-window load. We do not claim LLMs have human cognitive processes; we claim that notation properties reducing interpretive burden also reduce training-data ambiguity for statistical models.

\subsection{Research Hypotheses}

\begin{description}[noitemsep,style=nextline]
  \item[H1: Quality gradient.] LLM output quality decreases as formalization decreases. This is the natural prediction from the conceptual-modeling literature (formal representations constrain interpretation and should be easier for a reasoning agent). We test it and find it \emph{not} supported by the per-level automated metrics---an informative negative result, since the metrics turn out to be incommensurable and the formal level is in fact the hardest to get exactly right.
  \item[H2: Model--formalization interaction.] Different model architectures exhibit different performance profiles across formalization levels.
\end{description}

% ══════════════════════════════════════════════════════════════
\section{Method}
\label{sec:method}

\subsection{Domain: Northwind}

All three representations describe the Northwind food distribution domain: 11 entity types (Customer, Order, OrderLine, Product, Category, Supplier, Employee, Territory, Region, Shipper), 10 relationships, and business rules governing constraints such as positive quantities, valid discount ranges, and discontinued-product restrictions.

\subsection{High-Formal: OWL~2 Ontology + SPARQL}

The Northwind OWL~2 DL ontology contains 11 named classes, 20 object properties (with inverses), 45+ data properties, cardinality restrictions, disjointness assertions, and a populated ABox (5 customers, 8 orders, 15 order lines, 8 products, 3 employees, etc.). SWRL rules encode business constraints exceeding OWL~2 DL expressiveness.

\textbf{Tasks (100):} Given the ontology TBox and a natural-language competency question, generate the SPARQL query that answers it. Gold standard: hand-authored SPARQL queries verified by execution in Apache Jena against the populated ABox.

\textbf{Difficulty distribution:} Simple (20): class retrieval, property lookup. Medium (55): single-hop traversal, multi-hop join, aggregation, filtering. Hard (25): aggregation with HAVING, multi-hop aggregation, negation/optional patterns, constraint reasoning.

\subsection{Semi-Formal: EER Diagram + NL Annotations}

An EER diagram of the same domain, linearized as structured prose (entity/relationship/attribute declarations with participation and cardinality markers). Business rules appear as natural-language \texttt{Note:} annotations rather than formal OCL constraints. The annotations correspond directly to OWL axioms in the high-formal ontology, making the formalization gap traceable.

\textbf{Tasks (59):} Eleven categories: binary relationship constraint synthesis (6), ISA reasoning (5), participation inference (6), key identification (5), OCL synthesis from annotations (6), cardinality violation detection (5), weak entity recognition (5), ternary decomposition (5), derived attribute identification (5), redundancy detection (5), diagram consistency verification (6).

\textbf{OCL synthesis tasks} ask the model to formalize a natural-language annotation as an OCL invariant. Gold standard: valid OCL expressions verified by parse and typecheck.

\subsection{Low-Formal: NL Requirements Document}

A 7-paragraph requirements document describing the Northwind business in analyst prose, authored by informalizing each EER element. No diagrams, no formal notation.

\textbf{Tasks (31):} Nine categories: entity identification (4), relationship identification (4), implicit cardinality inference (4), constraint elicitation (3), ambiguity flagging (4), redundancy detection (3), scope boundary identification (3), anomaly detection (3), completeness assessment (3).

\textbf{Gold standard:} The Northwind EER diagram. Entity/relationship coverage is computed against the diagram's entity and relationship sets.

\subsection{Models}

Two open-weight instruction-tuned models: \textbf{Mistral-7B-Instruct-v0.3} (7.24B parameters) and \textbf{Llama-3.1-8B-Instruct} (8.03B parameters). Both run on NVIDIA RTX PRO 6000 Blackwell GPUs via HuggingFace Transformers with 4-bit quantization. Generation: $T\!=\!0.7$, $K\!=\!5$ repeated runs per task.

\subsection{Evaluation}

\begin{itemize}[noitemsep]
  \item \textbf{High-formal:} SPARQL correctness, executed in-process against the populated ABox. We report three nested metrics because they diverge sharply: \emph{runs} (executes without error), \emph{exact} (result set equals gold exactly), and \emph{answer} (projection-tolerant match: gold URIs and predicted labels are resolved to the same domain entities, so a query listing product \emph{names} counts as correct against gold listing product URIs). \emph{Answer} is our primary correctness metric; \emph{runs} is reported only to expose how misleading an executes-without-error metric is. Post-processing normalizes output (markdown stripping, prefix injection, variable-spacing fixes).
  \item \textbf{Semi-formal:} OCL presence (strict: \texttt{context X inv} pattern; loose: any \texttt{context X} pattern), parse success, and typecheck against a metamodel extracted programmatically from the full ontology (10 classes, 20 object properties, 45 data properties).
  \item \textbf{Low-formal:} Entity coverage against \emph{each task's own} gold entity set (extracted from the task's gold answer), not a single global entity set. This avoids penalizing excerpt tasks---which legitimately concern a subset of the domain---for omitting out-of-scope entities. Matching uses ontology-derived synonyms and is case-insensitive. Coverage is reported for the entity- and relationship-identification tasks (the extraction tasks); the remaining low-formal categories are reasoning/flagging tasks for which entity coverage is not the appropriate metric.
\end{itemize}

All three levels' automated metrics are incommensurable, so we additionally define a unified 0--3 expert rubric (Excellent/Adequate/Partial/Incorrect, with level-specific anchors) applied by two annotators to a stratified sample; this is the only instrument that supports a valid cross-level quality comparison (see the discussion of H1).

% ══════════════════════════════════════════════════════════════
\section{Results}
\label{sec:results}

\subsection{Dataset Summary}

Table~\ref{tab:dataset} summarizes the experimental data. All 190 tasks describe the same Northwind domain at three formalization levels.

\begin{table}[h]
\centering
\caption{Dataset summary. All levels share the Northwind domain.}
\label{tab:dataset}
\begin{tabular}{llrrrr}
\toprule
Level & Model & Tasks & $K$ & Records & Avg Len \\
\midrule
High-formal & Mistral   & 100 & 5 & 500 & 543 \\
High-formal & Llama     & 100 & 5 & 500 & 1{,}639 \\
Semi-formal & Mistral   &  59 & 5 & 295 & 740 \\
Semi-formal & Llama     &  59 & 5 & 295 & 3{,}691 \\
Low-formal  & Mistral   &  31 & 5 & 155 & 1{,}312 \\
Low-formal  & Llama     &  31 & 5 & 155 & 4{,}426 \\
\midrule
\multicolumn{2}{l}{Total} & 190 & & 1{,}900 & \\
\bottomrule
\end{tabular}
\end{table}

Response length increases monotonically with decreasing formalization for both models (Mistral: 543 $<$ 740 $<$ 1{,}312; Llama: 1{,}639 $<$ 3{,}691 $<$ 4{,}426), confirming the shift from terse formal queries to extended prose.

\subsection{High-Formal: SPARQL Execution (H1)}

Table~\ref{tab:sparql} presents SPARQL correctness under the three metrics of Section~\ref{sec:method}, by difficulty tier.

\begin{table}[h]
\centering
\caption{SPARQL correctness by difficulty tier (\% of $K\!=\!5$ runs over 100 tasks). \emph{runs} = executes without error; \emph{exact} = result set equals gold; \emph{answer} = projection-tolerant entity match (primary).}
\label{tab:sparql}
\setlength{\tabcolsep}{5pt}
\begin{tabular}{llrrrr}
\toprule
Model & Metric & Simple & Medium & Hard & \textbf{Overall} \\
\midrule
Mistral & runs   & 98\% & 70\% & 60\% & 73\% \\
        & exact  & 11\% & 6\%  & 11\% & 8\% \\
        & \textbf{answer} & \textbf{44\%} & \textbf{14\%} & \textbf{12\%} & \textbf{19\%} \\
\addlinespace
Llama   & runs   & 91\% & 76\% & 54\% & 74\% \\
        & exact  & 26\% & 5\%  & 13\% & 11\% \\
        & \textbf{answer} & \textbf{53\%} & \textbf{23\%} & \textbf{14\%} & \textbf{27\%} \\
\bottomrule
\end{tabular}
\end{table}

The three metrics diverge dramatically and this divergence is itself the finding. Both models generate queries that \emph{run} 73--74\% of the time, close to what a naive ``execution accuracy'' would report---but most queries that run are silently wrong: exact result-set match is only 8--11\%. The gap is largely a projection artifact (a query answering ``list all products'' by returning product \emph{names} rather than the gold's product URIs), which the projection-tolerant \emph{answer} metric corrects; even so, correctness is only 19\% (Mistral) and 27\% (Llama). All gold queries execute against the ABox, so the low scores reflect model error, not benchmark error. Correctness falls with task difficulty (Mistral 44\%$\to$12\%; Llama 53\%$\to$14\%), and Llama leads Mistral overall (27\% vs.\ 19\%), foreshadowing the model interaction (H2). The practical implication is immediate: formal SPARQL generation at the 7--8B scale requires execution verification, because ``the query ran'' is nearly uncorrelated with ``the query is right.''

\subsection{Semi-Formal: OCL Synthesis and EER Reasoning (H1, H2)}

\paragraph{OCL synthesis.}
Table~\ref{tab:ocl} presents OCL synthesis results under two evaluation criteria: \emph{strict} (requires the standard \texttt{context X inv:} pattern) and \emph{loose} (accepts any \texttt{context X} pattern with a constraint body).

\begin{table}[h]
\centering
\caption{OCL synthesis evaluation (30 outputs per model: 6 tasks $\times$ $K\!=\!5$).}
\label{tab:ocl}
\begin{tabular}{lcccc}
\toprule
& \multicolumn{2}{c}{Strict (\texttt{context X inv})} & \multicolumn{2}{c}{Loose (\texttt{context X})} \\
\cmidrule(lr){2-3}\cmidrule(lr){4-5}
& Presence & Parse & Presence & Typecheck \\
\midrule
Mistral & 8/30 (27\%) & 2/30 & 27/30 (90\%) & 13/30 \\
Llama   & 26/30 (87\%) & 26/30 & 28/30 (93\%) & 19/30 \\
\bottomrule
\end{tabular}
\end{table}

Under strict evaluation, Llama dominates (87\% vs.\ 27\%). Under loose evaluation, both models produce constraint expressions at similar rates (90\% vs.\ 93\%). The gap is \emph{syntactic fidelity}, not \emph{semantic comprehension}: Mistral writes \texttt{context~Order:~self.freight~>=~0} (missing the \texttt{inv} keyword), while Llama writes \texttt{context~Order~inv:~self.freight~>=~0} (correct OCL syntax). Both express the same constraint logic; Llama has higher surface-level OCL fidelity, likely from greater OCL exposure in its training corpus.

\paragraph{Non-OCL EER reasoning tasks} (53 tasks, 265 outputs per model) require reasoning about participation constraints, weak entities, cardinality violations, and diagram consistency without producing formal OCL. Both models perform well qualitatively: correctly identifying OrderLine as a weak entity, inferring cascade-delete implications from total participation, and detecting cardinality violations in sample data.

\subsection{Low-Formal: Entity Coverage (H1)}

Scored against each task's own gold entity set, entity/relationship coverage on the extraction tasks (entity identification and relationship identification) is high and nearly model-independent: 96\% for Mistral and 95\% for Llama (entity identification 96\%/94\%, relationship identification 96\%/97\%). A full-document extraction task, for example, elicits all ten domain entities from both models. This corrects a substantially lower figure obtained under a global-denominator scoring scheme, which penalized excerpt tasks for not mentioning entities outside their scope; under fair per-task scoring, informal entity extraction is one of the \emph{strongest}, not weakest, capabilities in the benchmark. (The remaining low-formal categories---ambiguity flagging, cardinality inference, redundancy detection---are reasoning tasks for which entity coverage is not the appropriate metric; they await the expert rubric.)

\subsection{Cross-Level Summary: The Formalization Gradient (H1)}

\begin{table}[h]
\centering
\caption{Cross-level primary metrics. All levels describe the same Northwind domain. The metrics are \emph{not commensurable} (see text); the table shows why a monotonic-gradient reading is unwarranted.}
\label{tab:summary}
\begin{tabular}{llrrr}
\toprule
Level & Primary Metric & Mistral & Llama & Combined \\
\midrule
High-formal & SPARQL answer correctness & 19\% & 27\% & 23\% \\
Semi-formal & OCL presence (loose) & 90\% & 93\% & 92\% \\
Low-formal & Entity coverage (per-task) & 96\% & 95\% & 95\% \\
\bottomrule
\end{tabular}
\end{table}

\textbf{H1 is not supported.} On defensible per-level metrics the values do not decrease with formalization---they \emph{increase} (23\% $\to$ 92\% $\to$ 95\%). We stress that these metrics are incommensurable: 23\% is projection-tolerant executable correctness, an unforgiving bar, while 92--95\% are presence and coverage rates that credit partial or merely-plausible output. The inversion therefore does not mean informal tasks are ``easier'' in an absolute sense; it means the formal level is the only one where we can \emph{measure} correctness stringently, and under stringent measurement it is the hardest. The lesson is methodological: a cross-level quality claim cannot be read off these proxies in either direction. Establishing whether a genuine formalization gradient exists requires a single instrument applied across levels---the unified 0--3 expert rubric (Section~\ref{sec:method}), whose stratified sample we have prepared and which is the natural next step. What the automated results \emph{do} support is a characterization of \emph{failure mode and verifiability} by formalization level, developed in Section~\ref{sec:discussion}.

\subsection{Model--Formalization Interaction (H2)}

Table~\ref{tab:interaction} reveals a striking interaction between model architecture and formalization level.

\begin{table}[h]
\centering
\caption{Model--formalization interaction. Bold indicates the stronger model at each level.}
\label{tab:interaction}
\begin{tabular}{lcc}
\toprule
Level & Mistral & Llama \\
\midrule
High-formal (SPARQL answer) & 19\% & \textbf{27\%} \\
Semi-formal (OCL strict) & 27\% & \textbf{87\%} \\
Semi-formal (OCL loose) & 90\% & \textbf{93\%} \\
Low-formal (entity coverage) & \textbf{96\%} & \textbf{95\%} \\
\bottomrule
\end{tabular}
\end{table}

At the formal level, Llama outperforms on answer correctness (27\% vs.\ 19\%). At the semi-formal level, Llama dominates strict OCL synthesis (87\% vs.\ 27\%) but the gap nearly vanishes under loose evaluation (93\% vs.\ 90\%). At the informal level, both models converge (95--96\% entity coverage). H2 is supported: model architecture interacts with formalization level, and---most tellingly at the semi-formal level---the interaction is concentrated in \emph{syntactic fidelity} (where Llama leads) rather than \emph{semantic comprehension} (where the models are similar). The interaction is strongest where formal syntax is demanded and disappears where output is free-form prose.

\subsection{Output Consistency ($K\!=\!5$)}

Neither model achieves meaningful lexical consistency. Average unique outputs per task is 4.6--5.0 across all conditions. Only Mistral on high-formal SQL achieves any perfect consistency (1/100 tasks). This replicates the near-zero consistency finding from prior work at $T\!=\!0.7$.

% ══════════════════════════════════════════════════════════════
\section{Discussion}
\label{sec:discussion}

\subsection{The Automation Boundary in the CM Pipeline}

The results identify a concrete automation boundary for 7--8B parameter LLMs operating over conceptual models:

\begin{itemize}[noitemsep]
  \item \textbf{SPARQL from OWL (19--27\% answer correctness):} Queries execute 73--74\% of the time but return the right entities only 19--27\% of the time; even simple competency questions reach only 44--53\%. ``Runs'' is nearly uninformative about correctness, so LLM-assisted SPARQL generation is viable \emph{only} behind execution-based verification.
  \item \textbf{OCL from EER (27--87\% strict, 90--93\% loose):} Both models understand the \emph{semantics} of constraints (90--93\% produce recognizable constraint expressions) but differ in \emph{syntactic fidelity} to the OCL standard. This suggests that LLM-assisted OCL synthesis is viable with a syntax-correction layer.
  \item \textbf{Entity extraction from NL (95--96\%):} At the informal level, scored against each task's own gold, both models recover nearly all in-scope entities. This is the most reliable capability in the benchmark---the opposite of what the formalization-gradient intuition predicts---though the metric credits mention rather than deeper modeling judgment, which the expert rubric will assess.
\end{itemize}

\subsection{Semantic Comprehension vs.\ Syntactic Fidelity}

The most theoretically interesting finding is the dissociation between semantic comprehension and syntactic fidelity. Under loose evaluation, both models produce correct constraint expressions at 90--93\% rates---they understand what the constraint \emph{means}. Under strict evaluation, Mistral drops to 27\% because it writes \texttt{context~Order:} instead of \texttt{context~Order~inv:}---a single missing keyword.

This has direct implications for CM tool design. A tool that requires syntactically perfect OCL will see Mistral as dramatically worse than Llama (27\% vs.\ 87\%). A tool that accepts near-OCL and applies syntactic normalization will see both models as nearly equivalent (90\% vs.\ 93\%). The choice of how strictly to evaluate formal output determines whether the automation boundary is perceived as a hard wall or a permeable membrane.

\subsection{SEQUAL-Grounded Interpretation}

The results refine, rather than simply confirm, the SEQUAL framework:
\begin{itemize}[noitemsep]
  \item High syntactic quality ($Q_s$) representations (OWL) enable the most \emph{stringent} evaluation---correctness is machine-decidable---but not the highest measured scores; decidability exposes that most syntactically valid queries are semantically wrong. Formal $Q_s$ buys verifiability, not accuracy.
  \item Semi-formal representations with mixed $Q_s$ (EER diagram structure + NL annotations) produce the model--formalization interaction: models trained on formal corpora (Llama) achieve higher syntactic fidelity, while the semantic content is similarly accessible to both.
  \item Low $Q_s$ representations (NL) eliminate the syntactic advantage entirely, and both models converge; high coverage scores here reflect forgiving pragmatic-quality ($Q_p$) proxies rather than verified semantic quality.
\end{itemize}

\subsection{Implications for CM Tool Builders}

\paragraph{SPARQL generation requires execution verification.} Answer correctness is only 19--27\% overall and 44--53\% even on simple competency questions, while 73--74\% of queries execute---so a deployment that trusts ``the query ran'' would accept mostly-wrong answers. Execution against the ABox with result inspection is mandatory, not optional, at the 7--8B scale.

\paragraph{OCL synthesis requires a normalization layer.} Both models understand constraint semantics (90\%+) but produce syntactically imprecise OCL. A tool should accept near-OCL output and apply syntactic correction (adding missing \texttt{inv} keywords, normalizing classifier names) rather than requiring perfect OCL generation.

\paragraph{NL-to-ER entity extraction is reliable; deeper modeling judgment is unverified.} At 95--96\% per-task entity coverage, identifying the entities and relationships present in informal requirements is well within reach at the 7--8B scale. This does not settle whether models make sound \emph{modeling} decisions (normalization, weak-entity recognition, constraint elicitation)---the reasoning-oriented low-formal tasks are exactly what the expert rubric targets. Coverage measures what is mentioned, not what is correctly modeled.

\paragraph{Model selection should consider formalization level.} The model--formalization interaction suggests that a CM tool processing formal artifacts should prefer Llama (or models with similar formal-notation training), while tools processing informal inputs may use either model interchangeably.

% ══════════════════════════════════════════════════════════════
\section{Limitations}
\label{sec:limitations}

\begin{enumerate}[noitemsep]
  \item \textbf{Automated metrics only.} The primary results rely on SPARQL execution, OCL parse/typecheck, and entity coverage. Expert annotation using a unified 0--3 rubric is in progress; until complete, semantic quality of non-OCL semi-formal outputs and the quality of low-formal reasoning are not formally assessed.
  \item \textbf{Two models at similar scale.} Both models are 7--8B parameters. Larger models (70B+) may alter the formalization gradient.
  \item \textbf{Northwind training-data exposure.} The Northwind domain is well-known. Model familiarity provides an upper bound on formal-level performance; the formalization gradient is robust even under this favorable condition, but novel domains may show steeper degradation.
  \item \textbf{Cross-level metric incommensurability.} SPARQL execution accuracy, OCL presence, and entity coverage are not on the same scale. The unified rubric (in progress) will enable direct cross-level comparison.
  \item \textbf{Single domain.} While the single-domain design eliminates the domain confound, it limits generalizability. Replication on additional domains is needed.
\end{enumerate}

% ══════════════════════════════════════════════════════════════
\section{Conclusion}
\label{sec:conclusion}

We presented the first benchmark for evaluating LLM reasoning over conceptual models at three formalization levels, using OWL, EER, and NL representations of a single domain. The results do \emph{not} support the intuitive formalization hypothesis (H1): on defensible per-level metrics, quality does not decrease as formalization decreases---formal SPARQL is in fact the hardest to get correct (19--27\% answer accuracy despite a 73--74\% execution rate), while informal entity extraction is the most reliable (95--96\%). Because these metrics are incommensurable, the right reading is not an inverted gradient but a caution: cross-level quality claims demand a single instrument, and we have prepared a unified 0--3 expert rubric for that purpose. What the benchmark establishes, with domain held constant, is that formalization determines \emph{verifiability and failure mode}---formal outputs fail detectably and so demand verification, informal outputs pass forgiving metrics---and that model architecture interacts with formalization level in syntactic fidelity but not semantic comprehension.

The model--formalization interaction (H2) reveals that different architectures have complementary strengths: Llama outperforms on formal SPARQL and strict OCL synthesis, while both models converge at the informal level. Critically, the dissociation between semantic comprehension (90\%+ for both models) and syntactic fidelity (27\% vs.\ 87\% strict OCL) identifies a specific, addressable gap: LLMs understand what constraints mean but cannot always produce them in the correct formal notation.

For the ER community, these results establish that the Mylopoulos tripartition is not merely a taxonomic convenience but a predictor of LLM performance over CM artifacts. The automation boundary is now empirically located: formal operations (SPARQL from OWL) are viable with verification; semi-formal operations (OCL from EER) require a syntax-correction layer; informal operations (entity extraction from NL) need human support. These findings are directly actionable for designers of AI-assisted conceptual modeling tools.

% ══════════════════════════════════════════════════════════════
\begin{thebibliography}{99}

\bibitem{mylopoulos1992conceptual}
Mylopoulos, J. (1992).
\newblock Conceptual Modelling and {T}elos.
\newblock In Loucopoulos, P. and Zicari, R., editors, \emph{Conceptual Modelling, Databases, and CASE}, pp.\ 49--68. Wiley.

\bibitem{wand1990ontological}
Wand, Y. and Weber, R. (1990).
\newblock An ontological model of an information system.
\newblock \emph{IEEE Trans.\ Software Engineering}, 16(11):1282--1292.

\bibitem{krogstie1995quality}
Krogstie, J., Lindland, O.I., and Sindre, G. (1995).
\newblock Towards a deeper understanding of quality in requirements engineering.
\newblock In \emph{CAiSE 1995}, LNCS 932, pp.\ 82--95. Springer.

\bibitem{krogstie2016quality}
Krogstie, J. (2016).
\newblock \emph{Quality in Business Process Modeling}. Springer.

\bibitem{moody2009physics}
Moody, D.L. (2009).
\newblock The ``Physics'' of Notations: Toward a Scientific Basis for Constructing Visual Notations in Software Engineering.
\newblock \emph{IEEE Trans.\ Software Engineering}, 35(6):756--779.

\bibitem{olive2007conceptual}
Oliv\'{e}, A. (2007).
\newblock \emph{Conceptual Modeling of Information Systems}. Springer.

\end{thebibliography}

\end{document}
